\section{The complete quotient census}
Theorems~\ref{structure} and~\ref{geometry} establish the universal marked group and its local sections before any subgroup enumeration is used. Put
\[
 H_0=\langle s_1,s_2,s_3\rangle,\qquad H_3=\langle s_0,s_1,s_2\rangle.
\]
The retained general endpoint theorem and the local quotient identifications in Theorem~\ref{geometry} give
\begin{equation}\label{criterion}
 N\text{ is semisparse}\quad\Longleftrightarrow\quad
 N^w\cap H_0H_3\subseteq\{1,\omega\}\quad(w\in W).
\end{equation}
The general quotient correspondence identifies these subgroup conjugacy classes with the quotient polytopes, and their inclusions with coverings. These general geometric inputs remain external; the particular group and its local marked identifications have been proved above. No assertion that an arbitrary polytopal orbit poset implies semisparseness is used.

The required program reconstructs the same direct-product marking as in Theorem~\ref{structure}, using the full Janko marking search and the source's literal twenty-point permutations. It checks all defining relators and both endpoint orders. Apply Lemma~\ref{avoid} to
\[
 B=H_0H_3\setminus\{1,\omega\}.
\]
The product has $1200$ elements, meeting $47$ of the $180$ element conjugacy classes of $W$ after the two permitted elements are removed. Every representative of every element class in each subgroup is tested against this full forbidden-class set. For a surviving subgroup, \eqref{criterion} also guarantees exactly the local alternatives required by Proposition~\ref{facets}. All vertex figures remain hemi-icosahedra; its $d$ facets are dodecahedra and its $h$ facets hemi-dodecahedra.

\begin{theorem}\label{census}
The group $W$ has $137$ conjugacy classes of semisparse subgroups, of which $26$ are maximal. The corresponding quotient polytopes comprise $65$ with the original facets and vertex figures, one with only hemi-dodecahedral facets, and $71$ with mixed facets. Thus $66$ are section regular. The eight rows removed in \cite{erratum} are exactly
\[
 13,\ 70,\ 78,\ 90,\ 91,\ 95,\ 111,\ 127.
\]
\end{theorem}
\begin{proof}
The required program uses GAP's algorithm enumerating every subgroup conjugacy class, yielding $1262$ representatives. It does not read the authors' old subgroup file, a semisparse list, or a list of quotient polytopes. Lemma~\ref{avoid} filters this exhaustive list to $137$ classes. Proposition~\ref{facets} gives $65$ rows with $h=0$, one with $d=0$, and $71$ with $d,h>0$. The two homogeneous facet cases are precisely section regular here, since the vertex figure is fixed and all lower sections are those of the two stated regular facets. The program separately enumerates maximal subgroups of every surviving representative. The resulting graph on ambient subgroup classes has $349$ distinct edges and, by Lemma~\ref{maxima}, $26$ maximal vertices.

For the source labels, all thirty generating sets in the original Table 1 are evaluated as literal words in this same marking. Exactly the four labelled $13,70,95,127$ fail the test; the other twenty-six are exactly the computed maximal classes. The four remaining deletions are identified from the printed inclusion rows as follows:
\begin{center}\begin{tabular}{cll}\toprule
Row&Printed subgroup&Identification in the literal parent\\\midrule
78&$C_{45}$&the unique subgroup class of order $45$ in row 13\\
111&$C_{15}$&the unique subgroup class of order $15$ in row 13\\
90&$C_{30}$&the unique cyclic subgroup class of order $30$ in row 70\\
91&$C_3\times D_{10}$&the order-$30$ nonabelian child of row 70\\
&&with normalizer quotient of order $12$\\\bottomrule\end{tabular}\end{center}
Here $D_{10}$ has order ten. The other nonabelian order-$30$ child has normalizer quotient of order $36$ and is not silently identified with row 91. For each of the eight rows the output gives actual permutations $x,q,b$ with
\[
 x\in N,\quad q\in W,\quad x^q=b\in H_0H_3\setminus\{1,\omega\}.
\]
These are explicit failures of the required universal conjugate condition.

Finally, the $137$ surviving historical rows, using all their printed maximal-subgroup links and facet counts, are matched bijectively to the computed $137$-vertex, $349$-edge graph. The thirty literal generating sets anchor the maximal rows; every surviving row is below a surviving maximum. The exact graph isomorphism, orders and both facet counts are checked, rather than only the grand total. All $137$ normalizer quotients also agree with the printed automorphism columns, after elementary direct-product identities. These include $D_{12}\cong C_2\times S_3$, $D_{20}\cong C_2\times D_{10}$, splitting coprime cyclic factors, and $\operatorname{GL}_2(4)\cong C_3\times A_5$: in the last identity the scalar center has order three, meets $\operatorname{SL}_2(4)\cong A_5$ trivially, and together these subgroups have the full order $180$. The complete original row set minus these survivors is exactly the displayed eight labels. The required finite output and the group/quotient premises above therefore establish all assertions.
\end{proof}

